Considera el pla $\pi\colon x+2y+2z-4=0$ i la recta
$r\colon\left\{\begin{aligned}x&=2\lambda\\y&=1-\lambda\\z&=2\end{aligned}\right.$

\begin{apartats}

\apartat[1,25]{0,75}
Comprova que $r$ és paral·lela a $\pi$ i calcula la distància entre tots dos.

\begin{solucio}
$\vec u=(2,-1,0)$ i $\vec n=(1,2,2)$: $\vec u\cdot\vec n=2-2+0=0$. El punt $(0,1,2)$ de $r$ no és de
$\pi$ ($0+2+4-4=2\ne0$): són paral·lels, i
\[
d(r,\pi)=\frac{|0+2+4-4|}{\sqrt{1+4+4}}=\frac23.
\]
\end{solucio}

\begin{tria}{a-distancia}
\itemtria{punts-a-un-terc}{1}{1,25}
Troba els punts de la recta $s\colon\left\{\begin{aligned}x&=t\\y&=t\\z&=1-t\end{aligned}\right.$ que
disten $\frac13$ del pla $\pi$.

\begin{solucio}
$\dfrac{|t+2t+2(1-t)-4|}{3}=\dfrac{|t-2|}{3}=\dfrac13$, és a dir, $|t-2|=1$: $t=3$ o $t=1$. Els punts són
$(3,3,-2)$ i $(1,1,0)$.
\end{solucio}

\itemtria{pla-que-conte-r}{1}{1,25}
Troba el pla que conté la recta $r$ i és paral·lel a $\pi$, i comprova que la distància entre els dos
plans és la distància entre $r$ i $\pi$.

\begin{solucio}
$x+2y+2z+D=0$ amb el punt $(0,1,2)$ de $r$: $2+4+D=0$, $D=-6$. El pla és $x+2y+2z-6=0$, i conté $r$
perquè hi conté un punt i la recta és paral·lela a $\pi$. Amb el punt $(0,0,2)$ de $\pi$:
$\dfrac{|4-6|}{3}=\dfrac23$, com a l'apartat anterior.
\end{solucio}
\end{tria}

\begin{nomesllarg}
\apartat{0,75}
Troba els valors de $k$ perquè el punt $(k,0,0)$ sigui a distància 1 del pla $\pi$.

\begin{solucio}
$\dfrac{|k-4|}{3}=1$, és a dir, $k-4=\pm3$: $k=7$ o $k=1$.
\end{solucio}
\end{nomesllarg}

\end{apartats}
